% GENERATED FILE. Edit canonical lessons, then run npm run generate:books.
% canonical-source-hash: fnv1a64:2ddbac4f78131ef5
% canonical-lessons: GE-C150-sparen, GE-C150-springen, GE-C150-sterben, GE-C150-stoeren, GE-C150-streiten

\chapter{To save, To jump, To die, To disturb, To argue}
\label{ch:ge-to-save-to-jump-to-die-to-disturb-to-argue}
\hlchaptermodality{150}

\begin{chapteropening}
\textbf{By the end of this chapter:} \emph{I can say to save, to jump, to die, to disturb and to argue.}

Complete the last lesson of chapter 150: I can say to save, to jump, to die, to disturb and to argue.
\end{chapteropening}

\section[sparen]{sparen — to save (money)}

\label{lesson:GE-C150-sparen}



\begin{quote}
Before the new one: say the German for to give, then the German for to shout.
\end{quote}

\subsection*{You'll want to know: sparen}
\textbf{sparen} — \textquotedblleft{}to save (money)\textquotedblright{}.

\begin{cousinweb}
\textbf{sparen} is inherited, from Old High German \emph{sparēn}. EtymDB 2.1 lists English \emph{spare} beside Dutch \emph{sparen} and Swedish \emph{spara}: one Germanic verb.
\end{cousinweb}

\begin{grammarlens}[title={What you've built}]
One more word for this chapter.
\end{grammarlens}

\subsection*{Your turn}
\begin{itemize}
\raggedright
  \item \emph{Say it:} \emph{sparen}
  \item \emph{Say it:} \emph{sparen}, once more
  \item \emph{Say it:} say \emph{sparen}
  \item \emph{Recall:} say the German for to give, then the German for to shout, then say \emph{sparen} again
\end{itemize}

\begin{tcolorbox}[breakable,colback=teal!4,colframe=teal!35!black,title={Before you move on}]
What does sparen mean? (\textquotedblleft{}To save (money)\textquotedblright{}.) Say it once more.
\end{tcolorbox}
\section[springen]{springen — to jump}

\label{lesson:GE-C150-springen}



\begin{quote}
Before the new one: say the German for to shout, then the German for to save.
\end{quote}

\subsection*{You'll want to know: springen}
\textbf{springen} — \textquotedblleft{}to jump\textquotedblright{}.

\begin{cousinweb}
\textbf{springen} $\leftarrow$ Old High German \emph{springan} $\leftarrow$ Proto-Germanic \emph{*springaną} — inherited, the same verb as English \emph{spring}, \textquotedblleft{}to leap\textquotedblright{}. (EtymDB 2.1)
\end{cousinweb}

\begin{grammarlens}[title={What you've built}]
One more word for this chapter.
\end{grammarlens}

\subsection*{Your turn}
\begin{itemize}
\raggedright
  \item \emph{Say it:} \emph{springen}
  \item \emph{Say it:} \emph{springen}, once more
  \item \emph{Say it:} say \emph{springen}
  \item \emph{Recall:} say the German for to shout, then the German for to save, then say \emph{springen} again
\end{itemize}

\begin{tcolorbox}[breakable,colback=teal!4,colframe=teal!35!black,title={Before you move on}]
What does springen mean? (\textquotedblleft{}To jump\textquotedblright{}.) Say it once more.
\end{tcolorbox}
\section[sterben]{sterben — to die}

\label{lesson:GE-C150-sterben}



\begin{quote}
Before the new one: say the German for to save, then the German for to jump.
\end{quote}

\subsection*{You'll want to know: sterben}
\textbf{sterben} — \textquotedblleft{}to die\textquotedblright{}.

\begin{cousinweb}
\textbf{sterben} $\leftarrow$ Old High German \emph{sterban} $\leftarrow$ Proto-Germanic \emph{*sterbaną}, \textquotedblleft{}to become stiff, to die\textquotedblright{} — inherited. (EtymDB 2.1) English \emph{starve} is the same verb, narrowed to dying of hunger.
\end{cousinweb}

\begin{grammarlens}[title={What you've built}]
One more word for this chapter.
\end{grammarlens}

\subsection*{Your turn}
\begin{itemize}
\raggedright
  \item \emph{Say it:} \emph{sterben}
  \item \emph{Say it:} \emph{sterben}, once more
  \item \emph{Say it:} say \emph{sterben}
  \item \emph{Recall:} say the German for to save, then the German for to jump, then say \emph{sterben} again
\end{itemize}

\begin{tcolorbox}[breakable,colback=teal!4,colframe=teal!35!black,title={Before you move on}]
What does sterben mean? (\textquotedblleft{}To die\textquotedblright{}.) Say it once more.
\end{tcolorbox}
\section[stören]{stören — to disturb}

\label{lesson:GE-C150-stoeren}



\begin{quote}
Before the new one: say the German for to jump, then the German for to die.
\end{quote}

\subsection*{You'll want to know: stören}
\textbf{stören} — \textquotedblleft{}to disturb\textquotedblright{}.

\begin{cousinweb}
\textbf{stören} $\leftarrow$ Middle High German \emph{stœren} $\leftarrow$ Old High German \emph{stōren}, \textquotedblleft{}to scatter, to destroy\textquotedblright{} — inherited. EtymDB 2.1 lists it beside English \emph{stir} and Dutch \emph{storen}.
\end{cousinweb}

\begin{grammarlens}[title={What you've built}]
One more word for this chapter.
\end{grammarlens}

\subsection*{Your turn}
\begin{itemize}
\raggedright
  \item \emph{Say it:} \emph{stören}
  \item \emph{Say it:} \emph{stören}, once more
  \item \emph{Say it:} say \emph{stören}
  \item \emph{Recall:} say the German for to jump, then the German for to die, then say \emph{stören} again
\end{itemize}

\begin{tcolorbox}[breakable,colback=teal!4,colframe=teal!35!black,title={Before you move on}]
What does stören mean? (\textquotedblleft{}To disturb\textquotedblright{}.) Say it once more.
\end{tcolorbox}
\section[streiten]{streiten — to argue}

\label{lesson:GE-C150-streiten}



\begin{quote}
Before the new one: say the German for to die, then the German for to disturb.
\end{quote}

\subsection*{You'll want to know: streiten}
\textbf{streiten} — \textquotedblleft{}to argue\textquotedblright{}.

\begin{cousinweb}
\textbf{streiten} $\leftarrow$ Old High German \emph{strītan}, \textquotedblleft{}to fight\textquotedblright{} $\leftarrow$ Proto-Germanic \emph{*strīdaną} — inherited; EtymDB 2.1 lists English \emph{stride} beside it.
\end{cousinweb}

\begin{grammarlens}[title={What you've built}]
One more word for this chapter.
\end{grammarlens}

\subsection*{Your turn}
\begin{itemize}
\raggedright
  \item \emph{Say it:} \emph{streiten}
  \item \emph{Say it:} \emph{streiten}, once more
  \item \emph{Say it:} say \emph{streiten}
  \item \emph{Recall:} say the German for to die, then the German for to disturb, then say \emph{streiten} again
\end{itemize}

\begin{tcolorbox}[breakable,colback=teal!4,colframe=teal!35!black,title={Before you move on}]
What does streiten mean? (\textquotedblleft{}To argue\textquotedblright{}.) Say it once more.
\end{tcolorbox}
